\section{Conclusões}
\label{sec-conclusoes}

O problema d-MBV é amplamente reconhecido como um desafio relevante na área de otimização. O objetivo central desta Iniciação Científica, em continuidade ao trabalho anterior do mesmo autor, foi aprimorar a busca de soluções para essa classe de problemas. As etapas previstas no projeto foram cumpridas e os resultados obtidos foram detalhados ao longo deste trabalho. O ponto de partida foi a tese desenvolvida por \cite{morenotese:2018}, que resultou na meta-heurística ILS. Com base nesse estudo, \cite{oliveira:2023} propôs um novo algoritmo construtivo, denominado R-BEP, que posteriormente foi aprimorado pelo mesmo autor do presente estudo \cite{silva2024ils}, originando a versão ILS+, uma meta-heurística mais robusta para o d-MBV. Cada uma dessas contribuições apresentou comparações próprias e evidenciou avanços nos resultados, estabelecendo um caminho contínuo de aprimoramento para a solução do problema.

Neste trabalho, foram desenvolvidas e analisadas as meta-heurísticas MULTI-START, GRASP-FIXO e GRASP-REATIVO, permitindo uma comparação geral com os resultados de pesquisas anteriores e buscando alcançar soluções ainda mais eficientes para o problema d-MBV. Os experimentos demonstraram que as três meta-heurísticas propostas superaram consistentemente os métodos de trabalhos anteriores, com destaque para o GRASP-REATIVO, que obteve os melhores resultados na maioria das instâncias. Esse desempenho evidencia que a estratégia de busca adaptativa adotada contribui de forma significativa para a qualidade das soluções, alcançando os resultados exatos em 314 de 525 instâncias. Além disso, ao comparar os valores obtidos entre o GRASP-REATIVO e os métodos ILS e ILS+, observou-se que o GRASP-REATIVO reduziu, em média, 1\% dos resultados em todos os conjuntos de instâncias. Considerando a complexidade intrínseca desse tipo de problema e seu caráter inteiro, esse desempenho representa uma melhoria relevante.

Como contribuição principal, este estudo amplia o conjunto de abordagens disponíveis para o problema d-MBV, demonstrando que diferentes estratégias de construção e intensificação podem levar a ganhos significativos na qualidade das soluções. Apesar dos bons resultados, algumas limitações ainda podem ser apontadas, como o elevado custo computacional em instâncias maiores e o tempo de execução mais longo quando o parâmetro de quantidade de iterações sem melhoria é aumentado. É importante ressaltar que, para atingir os melhores resultados, os custos de tempo foram expressivos, de modo que, dependendo do ambiente de aplicação, torna-se necessário avaliar o \textit{trade-off} entre qualidade e tempo de execução.

Entre possíveis trabalhos futuros, uma perspectiva promissora é a aplicação de técnicas de aprendizado de máquina para ajuste automático de parâmetros, bem como a exploração de variantes paralelas que permitam reduzir o tempo de execução em instâncias de grande porte. Outra linha de investigação seria a utilização do algoritmo construtivo desenvolvido neste trabalho combinado com diferentes estratégias de busca local, formando novas variantes de meta-heurísticas que possam aprimorar os resultados e reduzir os custos computacionais e temporais.